\documentclass[11pt]{article}
\usepackage{amsmath, amssymb}
\usepackage[margin=1in]{geometry}
\usepackage{fancyhdr}
\usepackage{hyperref}

\hypersetup{
    colorlinks=true,
    linkcolor=blue,      % Color of internal links (sections, pages)
    urlcolor=blue,        % Color of web URLs
    citecolor=green       % Color of citations
}

% ---- Fill these in ----
\newcommand{\courseName}{NEU 314}
\newcommand{\taName}{Keith Murray}
\newcommand{\sectionInfo}{Precept B01}
% ------------------------

\title{Precept 4: Warm Up Exercise}
\date{\today}

\pagestyle{fancy}
\fancyhf{}
\lhead{\courseName}
\chead{\sectionInfo}
\rhead{\today}
\cfoot{\thepage}

\begin{document}

\maketitle
\thispagestyle{fancy}

\section*{Problem}
Consider the following two matrices:
\[
R=\frac{1}{\sqrt{2}}\begin{bmatrix}
    1 & 1\\-1 & 1
\end{bmatrix},\qquad S=\begin{bmatrix}
    1 & 1\\0 & 1
\end{bmatrix}
\]
Let $\vec{e}_1=\begin{bmatrix}1\\0\end{bmatrix}$ and $\vec{e}_2=\begin{bmatrix}0\\1\end{bmatrix}$.
\begin{enumerate}
    \item Compute $R\vec{e}_2$ and $S\vec{e}_2$. How are they related? Now compute $R\vec{e}_1$ and $S\vec{e}_1$. What changed? Now compute $R^2\vec{e}_2$ and $S^2\vec{e}_2$. What would happen if we kept going (i.e. $R^n\vec{e}_2$ and $S^n\vec{e}_2$ for arbitrary $n$)?
    \item What are the eigenvectors of $R$? What are the eigenvectors of $S$? Try to answer geometrically first. For context, an eigenvector is a vector $\vec{x}$ such that $A\vec{x}=\lambda \vec{x}$
    \item Let's think about the singular value decomposition (SVD): $A=U\Sigma V^\top$. We will never learn how to compute this, so let's just think about it.
    \begin{enumerate}
        \item What would the SVD look like for $R$? \textbf{Hint:} What type of matrix is $R$?
        \item Imagine feeding every point on the unit circle into $S$. What shape comes out? What do the singular values $\sigma_1\geq\sigma_2$ correspond to in that picture?
        \item In words, describe what the SVD for $S$ would look like.
        \item Take a guess at what the determinant of $S$ is.
    \end{enumerate}
\end{enumerate}

\section*{Solution}
\subsection*{Question 1}
Computing the $R\vec{e}_2$ and $S\vec{e}_2$ gives us
\[
R\vec{e}_2=\frac{1}{\sqrt{2}}\begin{bmatrix}1\\1\end{bmatrix},\qquad S\vec{e}_2=\begin{bmatrix}1\\1\end{bmatrix}=\sqrt{2}\,R\vec{e}_2
\]
So $S\vec{e}_2$ points in exactly the same direction as $R\vec{e}_2$; it's just $\sqrt{2}$ times longer. If we only ever looked at $\vec{e}_2$, we might think that $S$ is a scaled version of $R$. But this is just a coincidence! Looking at $R\vec{e}_1$ and $S\vec{e}_1$, we have
\[
R\vec{e}_1=\frac{1}{\sqrt{2}}\begin{bmatrix}1\\-1\end{bmatrix},\qquad S\vec{e}_1=\begin{bmatrix}1\\0\end{bmatrix}=\vec{e}_1
\]
By now, it should be clear that $R$ is a rotation of $45^\circ$ clockwise. However, $S$ didn't rotate $\vec{e}_1$, it just outputted $\vec{e}_1$ again. What's going on with $S$?

Applying the matrices again to $\vec{e}_2$ we find
\[
R^2\vec{e}_2=\begin{bmatrix}1\\0\end{bmatrix},\qquad S^2\vec{e}_2=S\begin{bmatrix}1\\1\end{bmatrix}=\begin{bmatrix}2\\1\end{bmatrix}
\]
$R$ has now turned $\vec{e}_2$ by a total of $90^\circ$. But $S^2\vec{e}_2$ is only about $63^\circ$ from vertical, not $90^\circ$. In general,
\[
S^n=\begin{bmatrix}1 & n\\0 & 1\end{bmatrix},\qquad S^n\vec{e}_2=\begin{bmatrix}n\\1\end{bmatrix}
\]
If we keep applying $R$, the vector keeps going around the circle, with 8 applications taking us back to where we started $R^8=I$. If we keep applying $S$, the vector just gets longer and longer, and it leans closer and closer to the horizontal axis without ever crossing it. The matrix $S$ is what we call a \textit{shear} matrix.

\subsection*{Question 2}
Geometrically, an eigenvector is a vector that stays on its own line after being transformed. $R$ turns every vector by $45^\circ$, so no vector stays on its own line. Hence $R$ has \emph{no} real eigenvectors. (Algebraically, the eigenvalues of $R$ are complex, which we will ignore in this class.)

For $S$, we saw before that $\vec{e}_1$ did not get transformed. Hence,
\[
S\begin{bmatrix}1\\0\end{bmatrix}=1\cdot\begin{bmatrix}1\\0\end{bmatrix}
\]
and $\vec{e}_1$ is an eigenvector with eigenvalue $1$. Is there another one? 

For $\vec{v}=\begin{bmatrix}x & y\end{bmatrix}^\top$, the equation $S\vec{v}=\lambda\vec{v}$ gives us
\[
S\vec{v} = \begin{bmatrix}
    1 & 1\\0 & 1
\end{bmatrix}\begin{bmatrix}x\\y\end{bmatrix} = \begin{bmatrix}x+y\\y\end{bmatrix} = \begin{bmatrix}\lambda x\\ \lambda y\end{bmatrix} = \lambda\vec{v}
\]
where we now have the following system of equations
\[
x+y=\lambda x\qquad \text{and}\qquad y=\lambda y
\]
Either $y=0$ (giving $\vec{e}_1$ again), or $\lambda=1$, which forces $x+y=x$ and hence $y=0$, which contradicts $y\neq 0$. Therefore, $S$ has exactly one eigenvector (up to scaling). This is also why repeated application of $S$ pushed our vector toward the horizontal axis in Question 1.

\subsection*{Question 3}

\subsubsection*{Part (a)}
A rotation is already a rotation! So we can take $U=R$, $V^\top=I$, and $\Sigma=I$ to get
\[
R = RII = U\Sigma V^\top
\]
Notice that both singular values are $1$. Since rotations preserve the norm of input vectors, no direction is scaled.

\subsubsection*{Part (b)}
The unit circle gets mapped to a tilted ellipse. The singular values are the lengths of the \href{https://en.wikipedia.org/wiki/Semi-major_and_semi-minor_axes}{semi-axes of the ellipse}: 
\begin{itemize}
    \item $\sigma_1$ is how far the longest direction is stretched
    \item $\sigma_2$ is how much the shortest direction is squeezed
\end{itemize}
Also note that:
\begin{itemize}
    \item The columns of $U$ point along the axes of the ellipse
    \item For a rotation the ellipse is just the same circle, which is why $\sigma_1=\sigma_2=1$.
\end{itemize}

\subsubsection*{Part (c)}
As we did in \textbf{Part (b)}, imagine the input was a circle. Here is the procedure, reading $U\Sigma V^\top$ from right to left:
\begin{enumerate}
    \item Rotate the circle clockwise (this is $V^\top$)
    \item Stretch along the $x$ axis by $\sigma_1>1$ and squeeze along the $y$ axis by $\sigma_2<1$ to transform the circle into an ellipse (this is $\Sigma$)
    \item Rotate again, counterclockwise, to tilt the ellipse into its final position (this is $U$) 
\end{enumerate}
Even though the shear looks nothing like a rotation, it can be built out of two rotations and a stretch!

For the curious, the exact singular values are
\[
\sigma_1=\frac{1+\sqrt{5}}{2}\approx 1.618,\qquad \sigma_2=\frac{\sqrt{5}-1}{2}\approx 0.618
\]
where $\sigma_1$ is the golden ratio, and the SVD looks like
\[
S\approx\underbrace{\begin{bmatrix}0.851 & -0.526\\0.526 & 0.851\end{bmatrix}}_{U:\ \text{rotate } 31.7^\circ\text{ counterclockwise}}\underbrace{\begin{bmatrix}1.618 & 0\\0 & 0.618\end{bmatrix}}_{\Sigma}\underbrace{\begin{bmatrix}0.526 & 0.851\\-0.851 & 0.526\end{bmatrix}}_{V^\top:\ \text{rotate } 58.3^\circ\text{ clockwise}}
\]

\subsubsection*{Part (d)}
Notice that
\[
\sigma_1>1>\sigma_2
\]
which looks like 
\[
\sigma_1\cdot \sigma_2 \approx 1
\]
implying that 
\[
\det \lvert S\rvert = 1
\]
In fact, it is indeed the case that the determinant of $S$ is 1! We can show this via the exact singular values we noted above
\[
\frac{1+\sqrt{5}}{2}\frac{\sqrt{5}-1}{2} = \frac{-1 + \sqrt{5}-\sqrt{5}+ 5}{4}= \frac{4}{4}=1
\]
The determinant being equal to 1 implies that $S$ is volume preserving.


\end{document}
